\documentclass{beamer}[10]
\usepackage{pgf}
\usepackage{beamerthemesplit}

% Basado en: http://www.tjansson.dk/?p=82

\usepackage[utf8x]{inputenc}
\usepackage{moreverb}
\usepackage{hyperref}
\usepackage{xcolor}
\usepackage{listings}
\definecolor{mygreen}{rgb}{0,0.6,0}
\definecolor{mygray}{rgb}{0.5,0.5,0.5}
\definecolor{mymauve}{rgb}{0.58,0,0.82}
\definecolor{mygray}{rgb}{0.67, 0.73, 0.78}
\lstset{
	language=Python,
	basicstyle=\footnotesize,
	columns=fullflexible, % http://tex.stackexchange.com/questions/4911/phantom-spaces-in-listings-pdf
	upquote=true,
	breaklines=true,
	commentstyle=\color{mygreen},
	keywordstyle=\color{blue},
	stringstyle=\color{mymauve},
    extendedchars=true, 
    inputencoding=utf8x,
    literate={á}{{\'a}}1
             {é}{{\'e}}1
             {í}{{\'i}}1
             {ó}{{\'o}}1
             {ú}{{\'u}}1
             {ñ}{{\~n}}1
}

\definecolor{kugreen}{RGB}{149,122,143} %{50,93,61}
\definecolor{kugreenlys}{RGB}{132,158,139}
\definecolor{kugreenlyslys}{RGB}{173,190,177}
\definecolor{kugreenlyslyslys}{RGB}{214,223,216}
\setbeamercovered{transparent}
\mode<presentation>
\usetheme[numbers,totalnumber,compress,sidebarshades]{PaloAlto}

% Con número de página, o con footbar
%\setbeamertemplate{footline}[frame number]
\setbeamertemplate{footline}[frame]

  \usecolortheme[named=kugreen]{structure}
  \useinnertheme{circles}
  \usefonttheme[onlymath]{serif}
  \setbeamercovered{transparent}
  \setbeamertemplate{blocks}[rounded][shadow=true]
  \logo{\includegraphics[width=1.5cm]{imagenes/linti.pdf}}

% Ocultar barra de navegación
\beamertemplatenavigationsymbolsempty

\title{Programando con Python y Robots}
\author{Fernando López}
\institute{LINTI \\ Facultad de Informática \\ Universidad Nacional de la Plata
}
\date{24 de Octubre de 2013}


\begin{document}

\frame{\titlepage \vspace{-0.5cm}

}


\frame
{
\frametitle{Indice}
\tableofcontents[pausesection]
}


\begin{frame}{Actividad de entrega obligatoria}
	Para dar por aprobado el curso deberán entregar en grupos de 2 o 3 personas uno de los siguientes ejercicios:
	\begin{enumerate}[<+->]
		\item Planteo de 2 ejercicios para alumnos de escuelas secundarias.
		\item Planteo de 2 ejercicios para alumnos de nivel universitario.
		\item Planteo de modificaciones o ampliaciones posibles a la API del robot:
			\begin{itemize}
				\item Soporte para otros sensores/dispositivos.
				\item Añadir otros métodos o funcionalidades.
			\end{itemize}
	\end{enumerate}
	Enunciado en detalle: \url{https://github.com/Robots-Linti/curso-viedma#trabajo-final}
\end{frame}
\begin{frame}{Módulos}{Agrupando clases y funciones}
import x
import x as y
from x import y,z
from x import *
\verb+if __name__ == "__main__"+
\end{frame}
\begin{frame}{Paquetes}{Agrupando módulos}
\end{frame}
\begin{frame}{Interfaces gráficas}{Camino a un control remoto}
tk, Gtk3, Qt, WxWindows, otras...
Widgets
Window < Puede contener un solo elemento
Layouts > Box, Fixed, Grid, etc...
Button
Label
Scale > Adjustment
\end{frame}
\begin{frame}
Patterns:
Template pattern -> Joystick
Observer -> Scale
State -> Laberinto
\end{frame}
\frame
{
\frametitle{Direcciones}

\begin{block}{Direcciones de contacto}
\begin{center}
robots@linti.unlp.edu.ar \\
soportelihuen@linti.unlp.edu.ar
\end{center}
\end{block}

}


%~ \frame
%~ {
%~ \frametitle{Entorno de pre-boot del BIOS (3)}
	%~ Hoy en día sería deseable:
	%~ \begin{itemize}
		%~ \item Acceso al stack TCP, por ejemplo para gestores de 
		%~ arranque por red (en lugar de depender de la option ROM con 
		%~ soporte PXE que pudiera tener la NIC).
		%~ \item Uso de dispositivos como un mouse.
		%~ \item La posibilidad de ampliar o mejorar el soporte de 
		%~ hardware sin option ROMs.
		%~ \item Acceso a direccionamiento de 32 o 64 bits.
	%~ \end{itemize}
%~ }
%~ 
%~ \frame{
%~ 
%~ Efteråret 2004
%~ \subsection{Valg af projekt og vejleder}
%~ \frametitle{Tilfældighed}
%~ 
%~ \begin{block}{Kryptering med kaotiske kredsløb}
%~ \begin{columns}
%~ \column{.3\textwidth} \hspace{0.5cm}
%~ \includegraphics[width=0.7\textwidth]{imagenes/linti-f} 
%~ \column{.7\textwidth}
%~ \textit{Mogens Høgh Jensen}, NBI
%~ \end{columns}
%~ \end{block}
%~ 
%~ \begin{block}{Sandriller}
%~ \begin{columns}
%~ \column{.3\textwidth} \hspace{0.5cm}
%~ \includegraphics[width=0.7\textwidth]{imagenes/linti-f} 
%~ \column{.7\textwidth}
%~ \textit{Clive Ellegaard}, NBI og \emph{Tomas Bohr}, DTU
%~ \end{columns}
%~ \end{block}
%~ }





\end{document}
